% from ega4-18-fr.tex
\oldpage[IV-4]{166}

au nombre $n(y')$ (pour le morphisme $g$). Utilisant ensuite le fait que la propriété d'être universellement ouvert en un point se conserve par changement de base, on montre, comme dans (18.10.16, C)), que l'on peut se ramener au cas où $Y=\Spec(\sh{O}_{Y,y})$, avec $\sh{O}_{Y,y}$ strictement local. Alors, avec les notations de (18.10.16, D)), $X$ est somme des sous-préschémas ouverts $X_j=\Spec(\sh{O}_{X,x_j})$, où $x_j$ parcourt $g^{-1}(y)$; l'hypothèse que $g$ est ouvert en un de ces points $x_j$ entraîne que le sous-préschéma $X_j$ correspondant domine $Y$ (I, 10.3); on termine comme dans (18.10.16, D)).

\emph{On tire ensuite de (18.10.17.1) que les conclusions de (15.5.1) sont encore valables lorsqu'on ne suppose plus $Y$ localement noethérien, mais seulement que $f$ est séparé, quasi-fini et de présentation finie.}

En effet, pour prouver l'assertion (i) de (15.5.1), on remarque que, puisque $f$ est de présentation finie, l'ensemble $E$ des points $y'\in Y$ tels que $n(y')\geq n(y)$ est localement constructible (9.7.9); pour montrer que $y$ est intérieur à $E$, il suffit, en vertu de (I, 10.1), de voir que toute générisation $y'$ de $y$ appartient à $E$, ce qui n'est autre que (18.10.17.1).

Pour prouver l'assertion (ii) de (15.5.1), on peut, puisque $f$ est de présentation finie, utiliser (8.10.5, (xii)) appliqué suivant la méthode de (8.1.2, a)), et on peut donc déjà supposer que $Y=\Spec(\sh{O}_{Y,y})$. Appliquant ensuite (2.7.1, (vii)), on voit qu'on peut par le changement de base fidèlement plat $Y'\to Y$, où $Y'=\Spec({}^{hs}\!\sh{O}_{Y,y})$, supposer que l'anneau $\sh{O}_{Y,y}$ est strictement local. Appliquant (18.5.11, c)), on voit alors que si $x_j$ ($1\leq j\leq n$) sont les points de $f^{-1}(y)$, $X$ est somme de $n$ sous-préschémas ouverts $X_j=\Spec(\sh{O}_{X,x_j})$ qui sont finis sur $Y$, et d'un préschéma ouvert $X''$. Si on prouve que $X''=\varnothing$, on aura démontré que $f$ est fini, donc propre. Or, les corps $k(x_j)$ étant des extensions algébriques de $k(y)$, sont radicielles, et par suite $n(y)=n$; comme $f$ est ouvert en chacun des points $x_j$, $X_j$ domine $Y$ (I, 10.3), donc la restriction de $f$ à $X_j$, étant un morphisme fini, est surjective (II, 6.1.10). L'hypothèse que $z\mapsto n(z)$ est constant dans $Y$ entraîne donc que $X''\cap f^{-1}(y')=\varnothing$ pour tout $y'\in Y$, c'est-à-dire $X''=\varnothing$.

Enfin, la preuve de (iii) reproduit sans changement celle donnée dans (15.5.1).

\begin{env}[18.10.17.2]
\label{IV.18.10.17.2}
Par des méthodes analogues utilisant la hensélisation stricte, on peut, dans la plupart des résultats des \textsection\textsection~14 et 15, se débarrasser des hypothèses noethériennes.
\end{env}

\begin{lemma}[18.10.18]
\label{IV.18.10.18}
Soient $X$, $Y$ deux préschémas, $f:X\to Y$ un morphisme birationnel\footnote{Nous entendons par là la notion définie dans (6.15.4), mais où l'on ne suppose pas $X$ et $Y$ réduits.}, $x$ un point de $X$. Pour que $f$ soit un isomorphisme local au point $x$, il faut et il suffit que $f$ soit étale au point $x$. Pour que $f$ soit une immersion ouverte, il faut et il suffit que $f$ soit étale et séparé.
\end{lemma}

Les conditions énoncées étant trivialement nécessaires, tout revient à voir qu'elles sont suffisantes. Pour la première assertion, la question est locale sur $X$ et $Y$, donc on peut supposer $f$ étale, $X$ et $Y$ affines, $f$ séparé, et on est ramené à prouver la seconde assertion. En vertu de (17.9.1), il s'agit de prouver que $f$ est radiciel. Soit donc $y\in Y$ et prouvons que $f^{-1}(y)$ est radiciel sur $k(y)$. Le changement de base $\Spec(\sh{O}_{Y,y})\to Y$ ne changeant pas le fait que $f$ est étale, séparé et birationnel, on peut se borner au cas où $Y=\Spec(A)$, avec $A=\sh{O}_{Y,y}$.

